\subsection*{Algebraic Geometry and Commutative Algebra Notation}

\begingroup
\small

\noindent\emph{Status.} This part of the notes is still mostly scaffolded. The entries below reserve the notation promised by the section headings and keep it distinct from already-active algebra and analysis notation.

\notationgroup{Commutative Algebra}

\noindent\begin{tabularx}{\textwidth}{@{}p{0.38\textwidth}X@{}}
\notationentry{\operatorname{Spec}R,\ \mathfrak{p}\in\operatorname{Spec}R}{prime spectrum of a commutative ring and a prime ideal}
\notationentry{D(f),\ V(I)}{standard open set and closed set in the Zariski topology, when this part is developed}
\notationentry{S^{-1}R,\ R_{\mathfrak{p}},\ k(\mathfrak{p})}{localization at a multiplicative set, localization at a prime, and residue field}
\notationentry{\operatorname{rad}(I),\ \sqrt{I}}{radical of an ideal; use $\operatorname{rad}(\mathcal{V})$ for metric-vector-space radicals}
\end{tabularx}

\notationgroup{Affine and Projective Geometry}

\noindent\begin{tabularx}{\textwidth}{@{}p{0.38\textwidth}X@{}}
\notationentry{\mathbb{A}^n_k,\ \mathbb{P}^n_k}{affine and projective $n$-space over $k$; this is geometric notation, not probability}
\notationentry{I(X),\ Z(I),\ k[X]}{ideal of functions vanishing on $X$, zero set of an ideal, and coordinate ring}
\notationentry{f^\#,\ \varphi^\ast}{coordinate-ring pullback induced by a morphism; distinct from adjoints and form pullbacks by context}
\end{tabularx}

\endgroup
